% ---------------------------------------------------------------------
%  TÓM TẮT
% ---------------------------------------------------------------------
\begin{abstract}

Việc lựa chọn trang phục phù hợp thường đòi hỏi người dùng cân nhắc đồng
thời loại quần áo, màu sắc, phong cách và hoàn cảnh sử dụng. Đồ án tập trung
xây dựng hệ thống gợi ý phối đồ trên thiết bị di động nhằm hỗ trợ người dùng
nhận diện trang phục từ ảnh, quản lý tủ đồ cá nhân và lựa chọn các món đồ có
khả năng kết hợp với nhau. Hệ thống được định hướng gồm ứng dụng Flutter,
dịch vụ xử lý trí tuệ nhân tạo thông qua API và Firebase để lưu thông tin
người dùng, quần áo cùng các kết quả yêu thích.

Đầu vào của chức năng nhận diện là ảnh chứa một trang phục chính; đầu ra là
loại quần áo dự đoán và độ tin cậy. Trên cơ sở loại trang phục, màu sắc,
phong cách và dữ liệu trong tủ đồ, hệ thống dự kiến trả về tối đa ba phương
án phối đồ, đồng thời hạn chế kết quả trùng lặp. Đồ án cũng giới hạn chức
năng thử trang phục ảo ở kỹ thuật ghép ảnh hai chiều: ảnh quần áo có nền
trong suốt được đặt lên ảnh người dùng và có thể điều chỉnh vị trí, kích
thước, góc xoay. Cách tiếp cận này phù hợp với phạm vi học phần và không
được xem là mô hình sinh ảnh hay mô phỏng vật lý của vải.

Thành phần nhận diện sử dụng mô hình phát hiện đối tượng YOLO26n và giữ cố
định 21 lớp quần áo. Dữ liệu huấn luyện kết hợp 11.926 ảnh xuất từ Roboflow
với 52.712 ảnh DeepFashion In-shop; trong đó 52.711 ảnh DeepFashion được
chuyển thành công sang định dạng YOLO. Tập dữ liệu chuẩn bị cho phiên bản V3
gồm 59.463 ảnh huấn luyện, 6.902 ảnh xác thực và 6.620 ảnh kiểm thử. Ở lần
fine-tune cuối, mô hình đạt precision 55,18\%, recall 49,19\%, mAP@0,5 là
49,53\% và mAP@0,5:0,95 là 36,10\% trên tập xác thực. Thời gian phản hồi API
trên thiết bị Android thật chưa nằm trong phạm vi số liệu của lần chạy tích
hợp được sử dụng trong báo cáo.

\vspace{0.8em}
\noindent\textbf{Từ khoá:} gợi ý phối đồ, nhận diện quần áo, thị giác máy
tính, Flutter, thử trang phục ảo.

\end{abstract}
